\documentclass[11pt]{article}
\usepackage{microtype,amsmath,amssymb,amsthm,mathtools,booktabs,array}
\usepackage[margin=1.05in]{geometry}
\usepackage[colorlinks=true,linkcolor=blue!45!black,citecolor=blue!45!black,urlcolor=blue!45!black]{hyperref}
\usepackage{xcolor}
\hypersetup{pdftitle={The dimension-degree boundary for Gaussian moment vanishing},pdfsubject={Working mathematical manuscript}}
\newtheorem{theorem}{Theorem}[section]
\newtheorem{proposition}[theorem]{Proposition}
\newtheorem{lemma}[theorem]{Lemma}
\newtheorem{corollary}[theorem]{Corollary}
\theoremstyle{definition}\newtheorem{definition}[theorem]{Definition}
\theoremstyle{remark}\newtheorem{remark}[theorem]{Remark}
\newcommand{\E}{\mathbb E}
\newcommand{\C}{\mathbb C}
\newcommand{\tr}{\operatorname{tr}}
\newcommand{\wt}{\operatorname{wt}}
\newcommand{\G}{\mathsf G}
\DeclareMathOperator{\Span}{span}
\title{A Gaussian exclusion for hyperbolic suspensions in every degree}
\author{}\date{17 November 2026: internal written proof}
\begin{document}\maketitle
This note proves an all-degree family with an unbounded critical conic.
It lies outside the earlier gradient criterion for the phase $P$ itself.
The proof retains the full affine Gaussian,the complete ambient module,
all isolated degeneracies,and the original polynomial-marker filtration.
The classical tame rank comparison is credited to Pham. No novelty,
independent replay,formal verification or external review is claimed.

Let $q$ be a nondegenerate affine quadratic obtained from the standard
Gaussian by an invertible affine change of coordinates,and let its
normalized formal Gaussian functional be $\mathbb E$. For $a$ polynomial,
\[
 \Phi_a(t)=\sum_{m\ge0}\mathbb E[aP^m]\,t^m/m!.
\]
It is horizontal on the generic top quotient
\[
 M=\mathbb C(t)[x,y,w]/\sum_i(\partial_i+tP_i-q_i)\mathbb C(t)[x,y,w],
 \qquad\nabla_t=\partial_t+P.
\]
All changes of coordinates below retain this functional and the full
quadratic,including its means. No analytic period or image-kernel
calculation is required.

\begin{lemma}[A positive original seed]\label{lem:positiveformal}
Suppose a finite differential module $M$ has,after a finite cover
$r=t^{-1/N}$,a full frame of rank-one formal exponential factors.
If every zero-exponential horizontal functional takes a seed $a$ into
$r\mathbb C[[r]]$,the formal Gaussian functional cannot take $a$ to one.
\end{lemma}
\begin{proof}
If its value on $a$ is one,its restriction to the rational differential
submodule generated by $a$ takes values in $\mathbb C(t)$:the generator
values are $1,0,0,\ldots$. Thus it defines a horizontal quotient to the
trivial module. Formal extension is exact. After ramification the full
module is a direct sum of rank-one factors,so this quotient extends from
the seed submodule to the full module. A nonzero exponential factor
has no horizontal map to the trivial module over $\mathbb C((r))$:
the logarithmic derivative of a Laurent series cannot have its higher
polar derivative. On the zero part all horizontal maps are constant
linear combinations of the frame. Their seed values have positive order
and cannot sum to one.
\end{proof}

\begin{lemma}[Local isolated characters]\label{lem:isolatedcharacters}
For an exact local phase $\alpha+sz+y^{\mu+1}$,an analytic amplitude,
and an analytic volume unit,there are $\mu$ horizontal formal
characters. With $k_0=\mu+1$,the character indexed $j$, $1\le j<k_0$,
vanishes on a monomial $y^a s^b z^c$ unless $b=c$ and
$a=j-1+\ell k_0$, $\ell\ge0$. Otherwise it is
\begin{equation}\label{eq:ordinarycharacters}
 e^{\alpha t}t^{-1-j/k_0}
       (-1)^{b+\ell}b!\,(j/k_0)_\ell\,t^{-b-\ell}.
\end{equation}
\end{lemma}
\begin{proof}
At any given order only finitely many Taylor coefficients contribute.
The recurrences $I_{a+k_0}=-(a+1)I_a/(k_0t)$ and
$J_{b+1,b+1}=-(b+1)J_{b,b}/t$ annihilate all spatial exact forms.
The same recurrences verify that time differentiation inserts the
phase. The leading values are the jets $1,y,\ldots,y^{\mu-1}$,
multiplied by the analytic amplitude and volume unit. Their normalized
jet matrix is triangular with nonzero diagonal. Thus all multiplicities
are retained,without a Morse or distinct-critical-value hypothesis.
\end{proof}
\subsection{A critical conic in every polynomial degree}
\label{sec:hyperbolicsuspension}

The same fixed Laurent block handles a smooth critical conic in every
polynomial degree. We prove the full family,including a separate
rank proof and an explicit matrix establishing completeness.

\begin{theorem}\label{thm:hyperbolicsuspension}
Retain a full nondegenerate affine Gaussian quadratic
\[
 q=\tfrac12(Ax^2+By^2+mw^2+2Cxy+2Dxw+2Eyw)
                     +\ell_x x+\ell_y y+\ell_w w+q_0.
\]
Let $AB\ne0$ and
\[
 P=wxy+F(w)+k,\qquad F(0)=0,\quad F'(0)=h\ne0.
\]
Allow any polynomial degree $d_F=\deg F\ge1$, requiring $m\ne0$
when $d_F=1$. Then $P$ is not pure. Every original marked Gaussian
period is either zero or is not a finite exponential-polynomial.
The generic top twisted rank is $d_F+3$ for $d_F\ge2$ and five for
$d_F=1$. No Morse assumption on the Laurent phase obtained by
restricting $q$ to the conic, or on the isolated critical points of
$P$, is needed.
\end{theorem}

The nonzero quadratic poles at both ends of the conic are essential
to the rank-four construction below. The condition $AB\ne0$ states
this explicitly. Nondegeneracy of the full Gaussian alone would not
imply it. This theorem concerns complex polynomials and valued-field
formal frames; it does not assert an all-ring extension.

\begin{proof}
We first justify the full rank. Fix a generic $t\ne0$ and set
$S=tP-q$. Its gradient has a linear lower bound at complex spatial
infinity. To see this, suppose along an unbounded sequence
$\|\nabla S\|=o(R)$, where $R=\|(x,y,w)\|$. If $w$ stays
bounded, pass to a limit $w_0$ and normalize $(x,y)$ by $R$.
The first two gradient equations give
\[
 -AX+(tw_0-C)Y=0,\qquad (tw_0-C)X-BY=0.
\]
At least one of $X,Y$ is nonzero, and $AB\ne0$ makes both
nonzero. But the last gradient equation, divided by $R^2$, gives
$tXY=0$, a contradiction.

If $w$ is unbounded, the first two equations have a matrix with
inverse $O(1/|w|)$. Hence
$x,y=O(1)+o(R)/|w|$. First $R\asymp |w|$, then $x,y=O(1)$.
The last equation is therefore
\[
 tF'(w)-mw+O(1)=o(|w|).
\]
This is impossible for $d_F\ge3$, for $d_F=2$ with
$2t[w^2]F-m\ne0$, and for $d_F=1$ with $m\ne0$.
The same proof is uniform on compact subsets of these coefficient
conditions. Thus all phase-critical points stay in one ball during
any compact coefficient deformation inside that domain. They are
isolated, and their total Milnor number is constant by conservation
of local gradient-intersection multiplicity. Pham's tame Milnor
decomposition~\cite[Part~I, Sections~1--4]{Pham} identifies that
number with the full twisted rank.

For fixed $t$, the coefficient domain just used is the complement
of finitely many complex algebraic hypersurfaces, and is path
connected. We may compute the number at a generic diagonal
Gaussian with zero means, retaining $ABm\ne0$, and a generic $F$.
The critical equations there are
\[
 twy=Ax,\qquad twx=By,\qquad t(xy+F'(w))=mw.
\]
There are $d_F-1$ simple solutions with $x=y=0$ if $d_F\ge2$,
and one if $d_F=1$. The other solutions have
$w=\pm\sqrt{AB}/t$. At each of these two values, $x,y$ are
proportional and their nonzero product is prescribed by the last
equation; generically this gives two simple solutions. The counts
are $d_F+3$ and five, respectively. Uniform coercivity and the
connected deformation justify this specialization; a convenient
critical-point count by itself would not justify the rank for a
degenerate member.

We now construct a full formal frame. Set $r=t^{-1/2}$ and write
$F(w)=hw+bw^2+\sum_{j\ge3}f_jw^j$. Near the entire critical
conic $w=0$, $xy=-h$, use the exact coordinates
\[
 y=(rV-h)/x,\qquad w=rU.
\]
They give
\begin{align}
 tP-q&=k/r^2+UV+bU^2+\psi(x)+E_r(x,V,U),
                  &E_r\in r\C[x,x^{-1},V,U][r],\label{eq:conicphase}\\
 dx\,dy\,dw&=r^2\,\frac{dx}{x}\,dV\,dU,\label{eq:conicvolume}\\
 \psi(x)&=-\frac A2x^2-\frac{Bh^2}{2x^2}
                -\ell_x x+\frac{h\ell_y}{x}+Ch-q_0.\notag
\end{align}
All Gaussian cross terms and means are in $\psi$ or $E_r$.
The normal quadratic $UV+bU^2$ has Hessian determinant $-1$ for
every $b$. Its normalized formal Gaussian functional is well
defined and takes one to one.

The logarithmic Laurent quotient for $\psi$ has rank four, with
basis $1,x,x^2,x^3$. Here is a direct proof, including independence.
Its relation on $x^j$ is
\[
 jx^j-Ax^{j+2}-\ell_xx^{j+1}
             -h\ell_yx^{j-1}+Bh^2x^{j-2}=0.
\]
The top coefficient eliminates all exponents at least four without
creating a negative exponent; the bottom eliminates all negative
exponents without creating an exponent above three. If the image
of a nonzero Laurent polynomial under $x\partial_x+x\psi_x$
were supported in $[0,3]$, its highest exponent would be at most
one and its lowest at least two, which is impossible.

Let $\lambda_0,\ldots,\lambda_3$ be the dual basis on this
quotient, tensored with the normal Gaussian functional. For a
polynomial amplitude $a(x,y,w)$, define
\begin{equation}\label{eq:coniccharacters}
 L_i(a)=e^{k/r^2}r^2\lambda_i\!\left(
    e^{E_r}a(x,(rV-h)/x,rU)\right).
\end{equation}
Every coefficient is a finite Laurent-polynomial calculation, so
this is a well-defined formal series. Multiplication by $e^{E_r}$
conjugates the fixed twisted differential to the perturbed one.
Consequently these four functionals annihilate all spatial exact
forms. The moving volume is retained in the factor $r^2$.
For an explicit parameter check the coordinate velocity in the
original variables is
\[
 Z=\left(0,\frac{xy+h}{rx},\frac wr\right),
                   \qquad \operatorname{div}Z=2/r.
\]
Cartan's identity $\mathcal L_Z=d\iota_Z+\iota_Zd$ shows that
the pullback derivative differs from the original parameter
connection only by a spatial exact form and this volume
derivative. Thus \eqref{eq:coniccharacters} are horizontal for
the original full module. In particular, for every fixed
polynomial $a$,
\begin{equation}\label{eq:conicseedorder}
 L_i(a)\in e^{kt}r^2\C[[r]].
\end{equation}
This bound holds for all original markers, not just for the unit.

Suppose first $d_F\ge2$. The other critical points of $P$ are
$x=y=0$, $F'(w_p)=0$. Each $w_p$ is nonzero because $h\ne0$.
If it has multiplicity $\mu_p$, analytic coordinates give
\[
 P=k+F(w_p)+s z+y_0^{\mu_p+1}.
\]
Indeed absorb the nonvanishing factor $w$ into $x$ or $y$, then
use the one-variable power coordinate for $F$ at $w_p$.
The explicit characters \eqref{eq:ordinarycharacters}, including
their integration-by-parts and time recurrences, apply to this
exact local model. Their prefactors are
\[
 e^{(k+F(w_p))t}t^{-1-j/(\mu_p+1)},
                    \qquad 1\le j\le\mu_p.
\]
The Gaussian and coordinate Jacobian are analytic units in that
construction. These characters retain multiple roots and equal
critical values. Their total number is $\sum_p\mu_p=d_F-1$.

The $d_F+3$ characters just written are independent on polynomial
amplitudes. Put $B_0(w)=\prod_p(w-w_p)^{\mu_p}$. For the four
conic rows use $x^jB_0(w)$, $0\le j\le3$. Their leading
conic matrix is $B_0(0)I_4$. Their leading finite jets vanish.
For the remaining rows choose polynomials in $w$ with value zero
at $w=0$ and the separate Hermite jets of orders
$0,\ldots,\mu_p-1$ at the nonzero $w_p$. The Chinese remainder
theorem supplies them. Their conic values have an additional
positive power of $r$. Their normalized finite evaluation matrix
is the invertible local jet map, multiplied by analytic units.
After dividing each character column by its scalar prefactor,
the leading matrix is therefore block diagonal and invertible.
Since its size is the proved full rank, these are a full formal
frame. No uncomputed image kernel or infinity block remains.

For $d_F=1$ there is instead one further character. In the exact
coordinates
\[
 x=X/t,\qquad y=Y/t,\qquad w=(h/m)t+W,
\]
the phase is
\begin{align*}
 tP-q={}&\frac{h^2}{2m}t^2+(k-h\ell_w/m)t+H(X,Y,W)+O(t^{-1}),\\
 H={}&\frac hm XY-\frac hm(DX+EY)-\frac m2 W^2-\ell_w W-q_0.
\end{align*}
The correction is polynomial in $t^{-1},X,Y,W$, with positive
$t^{-1}$ order. The Hessian determinant of $H$ is $h^2/m\ne0$.
Formal Gaussian evaluation gives a horizontal character with
nonzero quadratic action and volume prefactor $t^{-2}$.
To check independence set $w_*=(h/m)t$. Use
$x^j(1-w/w_*)$, $0\le j\le3$, and $w/w_*$ as five polynomial
amplitudes with rational parameter coefficients. The former
have the original conic leading matrix and zero leading escape
values; the latter has zero leading conic values and nonzero
leading escape value. This gives the full rank-five frame.

In either degree case, after finite ramification the frame splits
the module into rank-one factors. Every zero-exponential
functional has positive order on every original polynomial
marker, by \eqref{eq:conicseedorder} and the displayed finite
prefactors. Lemma~\ref{lem:positiveformal} rules out a nonzero
constant original marked period. A constant-coefficient operator
reduces any nonzero finite exponential-polynomial to $c e^{\beta t}$;
it replaces the marker by another original polynomial in $P$.
Shifting $k$ by $-\beta$ and normalizing the marker gives the same
contradiction. This proves the theorem.
\end{proof}

\subsection{The cubic applications}
In the common ordinary-node cubic form
\[
 P=T(uv+h)+a u^3+d v^3+b u^2+c v^2+eu+fv+k,
\]
retain a full nondegenerate Gaussian with nonzero $T$ square coefficient.
Assume no leading node has both $F_2=0$ and zero Gaussian square
coefficient. Thus if $a=0,b=0$ the $u$ square coefficient is nonzero,
and if $d=0,c=0$ the $v$ square coefficient is nonzero. The critical
conic row is $h\ne0$, $a=b=c=d=e=f=0$. The two hyperbola rows
are $h=0$, $a=b=e=0$, $f\ne0$, and its symmetric counterpart.
These are the rows referred to in the following corollary.
\begin{corollary}\label{cor:ordinaryconics}
All the critical-conic and critical-hyperbola rows in the ordinary-node
table are not pure. Only its center-axis row remains open here.
\end{corollary}
\begin{proof}
For $h\ne0$ and $a=b=c=d=e=f=0$ the residual is
$T(uv+h)+k$. Apply the $d_F=1$ case with $w=T$, $x=u$, $y=v$.
The no-frame conditions give $ABm\ne0$.

For $h=0$, $a=b=e=0$, $f\ne0$, permute the coordinates and
write the residual as
\[
 v(Tu+d v^2+c v+f)+k.
\]
Take $w=v$, $x=T$, $y=u$, and
$F(w)=d w^3+c w^2+f w$. The quadratic coefficients along $T,u$
are nonzero: the original $T$ square coefficient is nonzero, and the $u$
node has $F_2=0$, so its Gaussian square coefficient is nonzero
by the no-frame condition. If this $F$ is linear, the remaining
$v$ node also has $F_2=0$, giving the extra $w$ square coefficient
required in degree one. The symmetric hyperbola row is identical
after swapping $u,v$. These are all the indicated rows.
\end{proof}
\begin{thebibliography}{9}
\bibitem{Pham}F.~Pham,\emph{La descente des cols par les onglets de
Lefschetz,avec vues sur Gauss--Manin},\emph{Ast\'erisque}
\textbf{130}(1985),11--47,Part~I,Sections~1--4;
\href{https://numdam.org/item/AST_1985__130__11_0/}{primary text}.
\end{thebibliography}\end{document}
